\documentclass[12pt]{article}

\usepackage{bbm}
\usepackage{fancyhdr}
\usepackage{titling}
\usepackage{enumitem}
\usepackage{amsmath}
\usepackage{dsfont}
\usepackage{amsfonts}
\usepackage{amssymb}
\usepackage{amsthm}
\usepackage{geometry}
\usepackage{tikz}
\usepackage{algorithm}
\usepackage{algpseudocode}
\geometry{verbose,tmargin=1in,bmargin=1in,lmargin=1in,rmargin=1in}
\renewcommand{\vec}[1]{\mathbf{#1}}
\newcommand{\dec}{\text{dec}}
\newcommand{\divi}{\mathbin{\textnormal{div}}}
\title{Homework 1}
\author{Satvik Prasad}
\date{}
\newtheorem{lemma}{Lemma}[section]
\fancypagestyle{plain}{%
  \fancyhf{}
  \fancyhead[L]{\theauthor}
  \fancyhead[R]{\thedate}
}
\pagestyle{plain}
\makeatletter
\def\@maketitle{%
  \newpage
  \null
  \vskip 1em%
  \begin{center}%
    \let \footnote \thanks
    {\LARGE \@title \par}%
    \vskip 1em%
  \end{center}%
  \par
\vskip 1em}
\makeatother
\usepackage{cmbright}

\newcommand{\mat}[1]{\mathbf{#1}}

% ---------- Probability theory ----------
% Core operators (optional argument-free; use as \E{X}, \Var{X}, \Cov{X}{Y})
\DeclareMathOperator{\EOp}{\mathbb{E}}
\DeclareMathOperator{\VarOp}{Var}
\DeclareMathOperator{\CovOp}{Cov}
\DeclareMathOperator{\CorrOp}{Corr}

\renewcommand{\P}[1]{\mathbb{P}\left(#1\right)}        % \P{A}
\newcommand{\E}[1]{\EOp\left[#1\right]}                % \E{X}
\newcommand{\Var}[1]{\VarOp\left(#1\right)}            % \Var{X}
\newcommand{\Cov}[2]{\CovOp\left(#1,\,#2\right)}       % \Cov{X}{Y}
\newcommand{\Corr}[2]{\CorrOp\left(#1,\,#2\right)}     % \Corr{X}{Y}
\newcommand{\SD}[1]{\operatorname{SD}\left(#1\right)}  % \SD{X}

% Subscripted versions, e.g. \Esub{\theta}{X}, \Psub{\theta}{A}
\newcommand{\Esub}[2]{\EOp_{#1}\left[#2\right]}
\newcommand{\Psub}[2]{\mathbb{P}_{#1}\left(#2\right)}

% Conditioning, indicators, independence
\newcommand{\given}{\,\middle|\,}                      % \E{X \given Y}
\newcommand{\ind}[1]{\mathbbm{1}\left\{#1\right\}}     % \ind{X > 0}
\newcommand{\indep}{\perp\!\!\!\perp}                  % X \indep Y
\newcommand{\iid}{\overset{\text{iid}}{\sim}}          % X_i \iid \Normal(0,1)

% Convergence
\newcommand{\convp}{\xrightarrow{\;p\;}}               % in probability
\newcommand{\convd}{\xrightarrow{\;d\;}}               % in distribution
\newcommand{\convas}{\xrightarrow{\text{a.s.}}}        % almost surely
\newcommand{\convLp}[1]{\xrightarrow{L^{#1}}}          % in L^p

% Probability spaces
\newcommand{\F}{\mathcal{F}}                           % sigma-algebra
\newcommand{\B}{\mathcal{B}}                           % Borel sigma-algebra
\newcommand{\pspace}{(\Omega, \F, \mathbb{P})}

% Common distributions
\newcommand{\Normal}{\mathcal{N}}
\newcommand{\Bern}{\operatorname{Bernoulli}}
\newcommand{\Bin}{\operatorname{Binomial}}
\newcommand{\Pois}{\operatorname{Poisson}}
\newcommand{\Unif}{\operatorname{Uniform}}
\newcommand{\Exp}{\operatorname{Exp}}
\newcommand{\Geom}{\operatorname{Geometric}}
\newcommand{\Gam}{\operatorname{Gamma}}
\newcommand{\Beta}{\operatorname{Beta}}

% Number sets
\newcommand{\R}{\mathbb{R}}
\newcommand{\N}{\mathbb{N}}
\newcommand{\Z}{\mathbb{Z}}


\begin{document}
\begin{enumerate}
  \item Let $h$ be a consistent heuristic, i.e, for all $v$ and
    neighbours of $v,$ $u,$ we have
    $$
    h(v) \leq c([v, u]) + h(u),
    $$
    where $c$ maps edges to their weights, and $h(G) = 0$ where $G$
    is our target vertex.

    Then, let us prove that $A*$ assigns the correct
    $\texttt{distance[v]} = \text{dist}(s, v)$ when $v$ is visited
    for the first time. Let us proceed by contradiction.

    Consider the first vertex $v$ assigned an incorrect distance
    $\texttt{distance[v]}.$ Let $u$ be the parent vertex of $v$ in
    the DAG produced by the edges taken by A*. Then, since there is
    no path from $s \to v$ with weight $< \text{dist}(s, v),$ we can conclude $$
    \texttt{distance[v]} > \text{dist}(s, v).
    $$

    Additionally, since $v$ is the first vertex with an incorrect
    assigned distance,
    $$
    \texttt{distance[u]} = \text{dist}(u, v).
    $$
    We take some edge $[u, v]$ to $v,$ meaning $\texttt{distance[v]}
    = \texttt{distance[u]} + c([u, v]),$ or equivalently,
    $$
    \texttt{distance[v]} = \text{dist}(s, u) + c([u, v]).
    $$

    Let $u'$ be the predecessor of $v$ in the correct minimum
    distance path from $s \to v,$ i.e., $\text{dist}(s, u') + c([u',
    v]) = \text{dist}(s, v).$ We have two cases:
    \begin{itemize}
      \item If $u'$ has already been visited, then
        $\texttt{distance[u']} = \text{dist}(s, u'),$ and the key
        $\text{dist}(s, u') + c([u', v]) + h(v)$ is in the heap. In
        this case, we require
        $$
        \texttt{distance[u]} + c([u, v]) + h(v) \leq
        \texttt{distance[u']} + c([u', v]) + h(v),
        $$
        since we took edge $[u, v]$ to $v$ before edge $[u', v].$
        Simplifying, we get
        $$
        \texttt{distance[u]} + c([u, v]) \leq \texttt{distance[u']} +
        c([u', v]).
        $$
        Since $u'$ is the predecessor on a minimum weight path from
        $s \to v,$ we must have
        $$
        \texttt{distance[v] = }\texttt{distance[u]} + c([u, v]) =
        \texttt{distance[u']} + c([u', v]) = \text{dist}(s, v),
        $$
        meaning $\texttt{distance[v]}$ is indeed set correctly,
        contradicting our assumption.

      \item Otherwise, $v$ is visited before $u'.$ Consider the
        earliest vertex $w$ along minimal $s \to u'$ that has been
        visited at the point $v$ is visited. Then,
        $$
        \texttt{distance[v]} + h(v) \leq \texttt{distance[w]} + h(w).
        $$
        By the consistence of $h,$ we further get
        $$
        h(w) \leq \text{dist}(w, u') + h(u'),
        $$
        (applying consistency along the the minimal path from $w \to
        u'$). Hence,
        $$
        \texttt{distance[v]} + h(v) \leq  \texttt{distance[w]} +
        \text{dist(w, u')} + h(u').
        $$
        Since $\texttt{distance[w]} = \text{dist(s, u)},$ we get
        $$
        \texttt{distance[v]} + h(v) \leq \text{dist}(s, u') + h(u')
        \leq \text{dist}(s, u')  + c[u', v] + h(v).
        $$
        Note that $u'$ is the predecessor of $v$ along the minimal path, meaning
        $$
        \text{dist}(s, u') + c[u', v] = \text{dist}(s, v),
        $$
        and thus
        $$
        \texttt{distance[v]} \leq \text{dist}(s, v),
        $$
        which is a contradiction.
    \end{itemize}

    In either case, we derive a contradiction, meaning $v$ cannot
    have an incorrect assigned distance. Hence, there is no first
    vertex $v$ with an incorrectly assigned distance, meaning all
    distances are assigned correctly if the vertex is visited.

  \item We use an approach similar to Djikstra's.

    \begin{algorithm}
      \begin{algorithmic}
        \Procedure{Solution}{$G$, $s$, $t$}
        \State $q := \texttt{max\_heap}()$
        \State $q \gets (\inf, [\texttt{nil}, s])$
        \State $\texttt{pred} := \texttt{map}()$

        \vspace{3mm}

        \While{$\neg \texttt{empty}(q)$}
        \State $\texttt{priority}, [u, v] \gets \texttt{pop}(q)$
        \If{$n \in \texttt{pred}$}
        \State \textbf{continue}
        \EndIf

        \State $\texttt{pred}[v] \gets u$

        \vspace{3mm}

        \For{$n \in G[v]$}
        \State $q \gets (\min(\texttt{priority}, c[v, n]), [v, n])$
        \EndFor
        \EndWhile

        \vspace{3mm}

        \State \Return $\texttt{as\_list}(\texttt{pred})$
        \EndProcedure
      \end{algorithmic}
    \end{algorithm}

    We prove that such an algorithm is correct, via contradiction.
    Assume that the path $p$ returned by our algorithm is incorrect.
    Let $\kappa(p)$ be the minimum edge weight of path $p,$ and let
    the path returned by our algorithm be $p,$ and a correct path be $p^\star.$

    Consider $e = [u, w],$ the edge in $p$ with weight $\kappa(p).$
    Since $c(e) = \kappa(p) < \kappa(p^\star),$ every edge in
    $p^\star$ has a higher priority, meaning we will never pop off
    edge $e$ from our priority queue until $p^\star$ has been
    traversed. However, this is a contradiction, since then our
    algorithm will assign $p^\star$ as our path from $s \to t.$

    Thus, the path $p$ returned must be one that maximizes $\kappa(p).$

    Since every edge is traversed once, and for every edge we push
    onto the priority queue of size $\mathcal O(|E|),$ the time
    complexity of this algorithm is $\mathcal O(|E| \log |E|).$

  \item We use the same approach as Djikstra's, and treat the keys in
    the priority queue as the minimum time required  to reach vertex $v.$

    \begin{algorithm}
      \caption{Delayed Djikstra's}
      \begin{algorithmic}
        \Procedure{DjikstraDelayed}{$G, W, S, R$}
        \State $q := \texttt{min\_heap}()$
        \State $q \gets (0, 0)$
        \State $\texttt{visited} := \texttt{set}()$

        \vspace{3mm}

        \While{$\neg \texttt{empty}(q)$}
        \State $\texttt{time}, v \gets \texttt{pop}(q)$

        \If{$v \in \texttt{visited}$}
        \State \textbf{continue}
        \EndIf

        \vspace{3mm}

        \State $\texttt{visited} \gets v$
        \If{$v \equiv n -1 $}
        \State \Return \texttt{time}
        \EndIf

        \vspace{3mm}

        \For{$w \in G[v]$}
        \If{$w \in \texttt{visited}$}
        \State \textbf{continue}
        \EndIf

        \vspace{3mm}

        \State $e \gets [w, n]$

        \If{$\texttt{time} + W[e] > R[e]$}
        \State \textbf{continue}
        \EndIf

        \vspace{3mm}

        \State $q \gets (\max(\texttt{time}, S[e]) + W_e, n)$
        \EndFor
        \EndWhile

        \vspace{3mm}

        \State \Return $\inf$
        \EndProcedure
      \end{algorithmic}
    \end{algorithm}

    Since this is just Djikstra's with a different priority that is
    still computed in $\mathcal O(1)$ time, the time complexity is
    $\mathcal O(|E| \log |E|).$

    We prove this algorithm is correct via contradiction. Assume that
    this algorithm returns the incorrect time required. Let $\tau
    (v)$ be the minimum time required to visit $v.$

    We prove that the keys in $\texttt{time}$ are always equal to
    $\tau(v).$ Consider the first vertex $w$ assigned an incorrect
    time $\texttt{time}[w] > \tau(w)$ (since no path can have a
    lesser time than the minimal time).

    Then, the predecessor of $w$ in the traversal, $u,$ was assigned
    a correct time. Consider the predecessor of $w$ in the optimal
    path, $u^\star.$

    \begin{itemize}
      \item If $u^\star$ was already visited, we know
        $$
        \max(\tau(u^\star), S[u^\star, w]) + W[u^\star, w] \geq
        \max(\tau(u), S[u, w]) + W[u, w],
        $$
        since edge $[u, w]$ was taken instead of edge $[u^\star, w].$

        Since $u^\star$ is the predecessor of $w$ in the minimal path, $$
        \max(\tau(u^\star), S[u^\star, w]) + W[u^\star, w] \leq
        \max(\tau(p), S[p, w]) + W[p, w]
        $$
        for all neighbors of $w,$ $p,$ since otherwise the path
        through $p$ is more minimal. Then,
        $$
        \max(\tau(u^\star), S[u^\star, w]) + W[u^\star, w] \leq
        \max(\tau(u), S[u, w]) + W[u, w].
        $$

        However, we chose $[u, w]$ instead of $[u^\star, w],$ and
        both sides of the expression are our priorities in the
        minimum heap, meaning
        $$
        \max(\tau(u^\star), S[u^\star, w]) + W[u^\star, w] \geq
        \max(\tau(u), S[u, w]) + W[u, w].
        $$

        Hence, since $\tau(w) = \max(\tau(u^\star), S[u^\star, w]) +
        W[u^\star, w]$ since because this edge is taken in the optimal path,
        \begin{align*}
          \tau(w) &= \max(\tau(u^\star), S[u^\star, w]) + W[u^\star,
          w] \\ &= \max(\tau(u), S[u, w]) + W[u, w] = \texttt{time}[w],
        \end{align*}
        which is a contradiction.

      \item Otherwise, $u^\star$ is visited after $w.$ Since
        $\texttt{time}[u^\star] \geq \tau(u^\star).$ At the time $w$
        is visited, consider the shortest path to $u^\star,$ and let
        the earliest vertex in the path not visited yet be $t.$ Then,
        since $w$ was popped off before $t$, and the predecessor of
        $t$, $t'$, has a correct $\texttt{time},$ we have
        $$
        \max(\tau(u), S[u, w]) + W[u, w] \leq \max(\tau(t'), S[t',
        t]) + W[t', t].
        $$

        Now, consider vertices following $t'$ in the optimal path to
        $u^\star.$ Since
        $$
        \max(\tau, S) + W \geq \tau,
        $$
        then $$
        \tau(u^\star) \geq \max(\tau(t'), S[t', t]) + W[t', t],
        $$
        meaning
        $$
        \texttt{time}[w]= \max(\tau(u), S[u, w]) + W[u, w] \leq
        \tau(u^\star) \leq \tau(w),
        $$
        since $\tau(w) = \max(\tau(u^\star), S[u^\star, w]) +
        W[u^\star, w].$  Once again, this is a contradiction, since
        $\texttt{time}[w] > \tau(w).$
    \end{itemize}

    In either case, we derive a contradiction, meaning there can be
    no such $w.$ Hence, every vertex is assigned a correct time,
    meaning $n - 1$ is assigned a correct time.

  \item When a new edge is added to the MST, the total weight
    increases by the weight of the path from the root to the new
    vertex incident by the edge.

    \begin{algorithm}
      \caption{Fancy MST}
      \begin{algorithmic}
        \Procedure{FancyMST}{$G$, $r$, $c$}
        \State $p,\ \texttt{mst} \gets \texttt{map}(),\ \texttt{set}()$
        \State $q := \texttt{min\_heap}()$
        \State $p[r] = 0$

        \vspace{3mm}

        \For{$n \in G[r]$}
        \State $q \gets (c[r, n], (r, n))$
        \EndFor

        \vspace{3mm}

        \While{$\neg \texttt{empty}(q)$}
        \State $\_, (u, v) \gets \texttt{pop}(q)$

        \vspace{3mm}

        \If{$v \in p$}
        \State \textbf{continue}
        \EndIf

        \vspace{3mm}

        \State $\texttt{mst} \gets (u, v)$
        \State $p[v] = p[u] + c[u, v].$

        \vspace{3mm}

        \For{$n \in G[v]$}
        \State $\Delta c \gets p[v] + c[v, n]$

        \State $q \gets (\Delta c, (v, n))$
        \EndFor
        \EndWhile

        \vspace{3mm}

        \State \Return \texttt{mst}
        \EndProcedure
      \end{algorithmic}
    \end{algorithm}

    We traverse every edge once, and push onto a heap of size
    $\mathcal O(|E|),$ yielding a complexity of $\mathcal O(|E| \log |E|).$

    We prove correctness via contradiction. Consider the first
    iteration where the set $\texttt{mst}$ was not the minimum fancy
    spanning tree of the visited set (keys of $\texttt{p}$). Let the
    minimum fancy spanning tree of vertex set $V$ be $\mu(V).$

    Then, our first invalid iteration has vertex set $V \cup \{v\}$,
    where $V$ is the vertex set of the previous iteration. Consider
    MFST $\mu(V \cup \{v\}),$ with some edge $e$ incident on $v$ from
    component $V.$ We know $$
    c[\mu(V \cup \{v\})] < c[\texttt{mst}]
    $$
    during this current iteration.

    Additionally, $\mu(V \cup \{v\}) = \mu(V) \cup e$ for some edge
    $e$ with incident vertex $v,$ and outgoing from some vertex $u
    \in V$. Since $\mu(V \cup \{v\})$ is minimal, and $$c[\mu(V \cup
    \{v\})] = c[\mu(V)] + c[p],$$ where $p$ is the path from $r$
    through the MST $\mu(V)$ through $e$ into $v,$ the chosen $[u,
    v]$ must minimize
    $$
    c[r \to \ldots \to u \to v].
    $$
    Since on this iteration we choose such an $e,$ we can conclude
    $c[\texttt{mst}] = c[\mu(V \cup \{v\})],$ which is a contradiction.

    Thus, there can be no such invalid iteration, and $\texttt{mst}$
    represents the MFST of the entire graph $G$ by the final
    iteration of our algorithm.
\end{enumerate}
\end{document}
